\documentclass[11pt]{article}

% --- Folienformat wie Game Concept Document (16:9, 960x540pt) ---
\usepackage[paperwidth=960bp, paperheight=540bp, margin=0pt]{geometry}
\usepackage[ngerman]{babel}
\usepackage[utf8]{inputenc}
\usepackage[T1]{fontenc}
\usepackage[scaled]{helvet}
\renewcommand{\familydefault}{\sfdefault}
\usepackage[autostyle=true,german=quotes]{csquotes}
\usepackage{amsmath,amssymb}
\usepackage{graphicx}
\usepackage{xcolor}
\usepackage{tikz}
\usetikzlibrary{arrows.meta, calc, shapes.geometric}
\usepackage{microtype}
\usepackage{hyperref}
\urlstyle{same}

\definecolor{hkared}{RGB}{198, 42, 18}
\definecolor{hkagray}{RGB}{100, 100, 100}

% --- Farben & Stile der Themenlandschaftskarte (aus Themenkarte_Agora.tex) ---
\definecolor{clBlue}{RGB}{45, 115, 195}
\definecolor{clGold}{RGB}{215, 165, 30}
\definecolor{clGreen}{RGB}{42, 155, 85}
\definecolor{clPurple}{RGB}{130, 95, 195}
\definecolor{clRed}{RGB}{215, 70, 65}
\tikzset{
  pfeil/.style     = {-{Stealth[length=2.2mm, width=1.6mm]}, line width=0.8pt, draw=black!80},
  looppfeil/.style = {-{Stealth[length=2.3mm, width=1.65mm]}, line width=1.0pt, draw=#1!85!black},
  txt/.style       = {font=\tiny\bfseries, fill=white, inner sep=1.3pt, rounded corners=1.5pt,
                      align=center, text=black!85},
  agoraZentrum/.style = {ellipse, draw=clGold!80!black, line width=2.2pt, fill=clGold!20,
                      text=black!85, minimum width=4.6cm, minimum height=2.25cm,
                      align=center, inner sep=3pt},
  bknoten/.style   = {draw=#1, line width=1.3pt, rounded corners=6pt, fill=#1!10,
                      text width=3.9cm, minimum height=1.25cm, align=center, inner sep=3pt,
                      font=\scriptsize},
}
% --- Schalter & Makros für die Agora-Systemkarte (agora_systemkarte.tex) ---
\newif\ifskKern
\newcommand{\pTag}{\textbf{[+]}}
\newcommand{\mTag}{\textbf{[--]}}
\newcommand{\pmTag}{\textbf{[$\pm$]}}
\newcommand{\zg}[1]{\\[0pt]{\tiny\itshape\color{black!62}#1}}

\hypersetup{
    colorlinks = true, linkcolor = hkared, urlcolor = hkared,
    pdftitle   = {Zeitgeist Handel},
    pdfauthor  = {Erlind Sejdiu},
    pdfsubject = {Game Design, Game Concept Document}
}

\pagestyle{empty}
\setlength{\parindent}{0pt}
\setlength{\parskip}{0.75em}

% --- Inhaltsfolie: Titelbalken links, Logo rechts unten, Seitenzahl ---
\newcounter{folie}
\setcounter{folie}{1}
\newcommand{\folie}[2]{%
  \stepcounter{folie}%
  \newpage
  \begin{tikzpicture}[remember picture, overlay]
    % Akzentbalken + Titel
    \fill[hkared] ([xshift=38bp, yshift=-18bp]current page.north west) rectangle ++(9bp, -64bp);
    \node[anchor=west, font=\fontsize{34}{40}\selectfont\bfseries, text=hkared]
      at ([xshift=58bp, yshift=-52bp]current page.north west) {#1};
    % Schatten-Linie unter Kopf
    \shade[top color=black!12, bottom color=white]
      ([yshift=-100bp]current page.north west) rectangle ([yshift=-108bp]current page.north east);
    % Logo
    \node[anchor=south east, inner sep=0pt]
      at ([xshift=-22bp, yshift=20bp]current page.south east)
      {\includegraphics[height=120bp, trim=520 140 580 140, clip]{HKA.jpg}};
    % Seitenzahl
    \node[anchor=south west, font=\scriptsize, text=hkagray]
      at ([xshift=60bp, yshift=24bp]current page.south west) {\thefolie};
    % Textblock
    \node[anchor=north west, inner sep=0pt, text width=760bp, align=justify,
          font=\fontsize{11.5}{15.5}\selectfont]
      at ([xshift=60bp, yshift=-130bp]current page.north west)
      {\parbox{760bp}{\setlength{\parskip}{0.8em}\justifying #2}};
  \end{tikzpicture}%
}
\usepackage{ragged2e}

% --- Kartenfolie: gleicher Kopf wie \folie, Inhalt als TikZ-Grafik ---
\newcommand{\kartenfolie}[2]{%
  \stepcounter{folie}%
  \newpage
  \begin{tikzpicture}[remember picture, overlay]
    \fill[hkared] ([xshift=38bp, yshift=-18bp]current page.north west) rectangle ++(9bp, -64bp);
    \node[anchor=west, font=\fontsize{34}{40}\selectfont\bfseries, text=hkared]
      at ([xshift=58bp, yshift=-52bp]current page.north west) {#1};
    \shade[top color=black!12, bottom color=white]
      ([yshift=-100bp]current page.north west) rectangle ([yshift=-108bp]current page.north east);
    \node[anchor=south east, inner sep=0pt]
      at ([xshift=-22bp, yshift=20bp]current page.south east)
      {\includegraphics[height=120bp, trim=520 140 580 140, clip]{HKA.jpg}};
    \node[anchor=south west, font=\scriptsize, text=hkagray]
      at ([xshift=60bp, yshift=24bp]current page.south west) {\thefolie};
  \end{tikzpicture}%
  {\hyphenpenalty=10000 \exhyphenpenalty=10000
  \begin{tikzpicture}[remember picture, overlay,
      shift={($(current page.north west)+(1.45cm,-0.45cm)$)},
      scale=0.88, every node/.append style={transform shape}]
    #2
  \end{tikzpicture}}%
}

\begin{document}

% ================= Titelfolie =================
\begin{tikzpicture}[remember picture, overlay]
  \fill[hkared] (current page.south west) rectangle (current page.north east);
  \node[anchor=north west, align=left, text=white, font=\scriptsize]
    at ([xshift=30bp, yshift=-28bp]current page.north west)
    {\textbf{Hochschule Karlsruhe}\\University of\\Applied Sciences\\[4pt]
     Fakultät für\\\textbf{Informatik und}\\\textbf{Wirtschaftsinformatik}};
  \node[anchor=west, text=white, font=\fontsize{24}{28}\selectfont]
    at ([xshift=66bp, yshift=-196bp]current page.north west) {Game Concept Document};
  \draw[white, line width=0.8pt]
    ([xshift=68bp, yshift=-238bp]current page.north west) -- ++(300bp, 0);
  \node[anchor=west, text=white, font=\fontsize{24}{28}\selectfont]
    at ([xshift=66bp, yshift=-280bp]current page.north west) {Zeitgeist Handel};
  \node[anchor=west, text=white, font=\large\bfseries]
    at ([xshift=56bp, yshift=-410bp]current page.north west) {Game Design | INFM | WS 2026/27};
  \node[anchor=west, text=white, font=\large]
    at ([xshift=56bp, yshift=-456bp]current page.north west) {Erlind Sejdiu};
  % Weiße Logospalte rechts
  \fill[white] ([xshift=-170bp]current page.north east) rectangle (current page.south east);
  \node[anchor=north, inner sep=0pt]
    at ([xshift=-85bp, yshift=-30bp]current page.north east)
    {\includegraphics[height=170bp, trim=520 140 580 140, clip]{HKA.jpg}};
\end{tikzpicture}
%
%% ================= Zeitgeist 1: Was ist Handel? =================
%\folie{Zeitgeist}{%
%Ein Markt ist das Spiegelbild der Gesellschaft. Die antike \textbf{Agora} war politischer Treffpunkt, ein sozialer mittelpunkt in diesem aller Art von Menschen zusammenkommen. 
%Seien es Kaufleute, Gaukler oder auch Richter, ein lebensnotwendiger Ort, damit man Lebensmittel erhält und ein Ort an dem politsche Systeme durch Institionen einfluss nehmen können. 
%Damit das alles funktioniert, braucht es ein zugrunde liegendes System. 
%Man muss dem System, was geben damit auch was zurückbekommt. Damit man in diesem System zurecht kommt, muss man die Regeln kennen und entsprechend nach diesen spielen. 
%
%Heute ist der Marktplatz fast vollständig aufs Smartphone gewandert. Bei Amazon, Temu oder Shein lässt sich beinahe alles mit einem
%Klick bestellen. Allein in Deutschland wurden 2024 täglich rund 400.000 Pakete von Temu und Shein zugestellt. Laut Bundesbank wurden
%2025 erstmals mehr als die Hälfte aller Einkäufe bargeldlos bezahlt, nämlich 55\,\%. Jeder dieser Käufe hinterlässt eine Datenspur.
%Der Soziale Treffpunkt ist ebenso ins Internet ausgewandert. Stats sich auf einem Marktplatz zu treffen, trifft man sich nur Digital. 
%Wer im Internet wen trifft, kann jetzt von Konzernen gesteuert werden
%
%
%Letzlich existiert die Agora immernoch es ist ein Treffpunkt, aber die Regel, um in diesem System erfolgreich zu sein ändern sich.
%
%}
%
%% ================= Zeitgeist 2: Gefühle & Ambivalenz =================
%\folie{Zeitgeist}{%
%Beim Thema Handel sind die Gefühle gemischt. Kaufen und Verkaufen macht vielen Menschen Spaß. Wer ein gutes Geschäft macht, ein seltenes
%Stück ergattert oder billig einkauft und teuer weiterverkauft, fühlt sich clever und erfolgreich. Plattformen nutzen das gezielt und
%greifen dafür zu Mitteln aus dem Game Design. Einkaufsapps zeigen Glücksräder und Countdowns, bei limitierten Drops konkurrieren Kunden
%um wenige Paar Sneaker, und auf StockX werden Turnschuhe mit Kurscharts gehandelt wie Aktien. So wird Handel zum Spiel, zu dem auch
%Nervenkitzel, Status und die Angst gehören, etwas zu verpassen (FOMO).
%
%Zugleich wächst das Unbehagen. Viele wollen fair und nachhaltig einkaufen und greifen bei steigenden Preisen dann doch zur Billigware,
%was ein schlechtes Gewissen hinterlässt. Der Staat reagiert mit Lieferkettengesetzen wie dem deutschen LkSG und der europäischen
%Richtlinie CSDDD, die Ausbeutung entlang der Lieferkette verhindern sollen. Dazu kommt das Gefühl, selbst nur Spielfigur zu sein.
%Algorithmen kennen die eigenen Vorlieben, Preise passen sich individuell an, und Bots kaufen begehrte Ware in Millisekunden weg, bevor
%ein Mensch überhaupt klicken kann.
%
%Handel verspricht Freiheit, Wohlstand und unbegrenzte Auswahl. Er macht aber auch abhängig von fernen Lieferketten, von großen
%Plattformen und vom eigenen Konsumdrang. Viele kennen Gedanken wie „Eigentlich brauche ich das nicht, aber es ist gerade so günstig“
%oder „Wer verdient hier eigentlich an mir?“
%
%Der Markt wirkt wie ein Ort, an dem jeder mit Geschick gewinnen kann. Wer am Ende tatsächlich profitiert, hängt in der Realität aber oft
%davon ab, wer mehr weiß, wer die Infrastruktur kontrolliert und wem die anderen vertrauen.
%}
%
%
%\folie{Zeitgeist}{%
%Ein Markt ist ein Spiegelbild der Gesellschaft. Die antike \textbf{Agora} war nicht nur ein Handelsplatz, sondern auch politischer
%Treffpunkt und sozialer Mittelpunkt. Hier kamen alle zusammen, Kaufleute, Gaukler und Richter. Der Markt versorgte die Menschen mit Lebensmitteln, und gleichzeitig nahmen Politik und Institutionen über ihn Einfluss auf
%das Zusammenleben.
%
%Damit das funktioniert, braucht es ein System, das auf Geben und Nehmen beruht. Wer etwas bekommen will, muss etwas einbringen, sei es
%Ware, Arbeit, Geld oder Vertrauen. Wer in diesem System zurechtkommen will, muss seine Regeln kennen und nach ihnen spielen.
%
%Heute ist der Marktplatz fast vollständig aufs Smartphone gewandert. Bei Amazon, Temu oder Shein lässt sich beinahe alles mit einem
%Klick bestellen. Allein in Deutschland wurden 2024 täglich rund 400.000 Pakete von Temu und Shein zugestellt. Laut Bundesbank wurden
%2025 erstmals mehr als die Hälfte aller Einkäufe bargeldlos bezahlt, nämlich 55\,\%. Jeder dieser Käufe hinterlässt eine Datenspur,
%und so bezahlt man neben Geld auch mit Daten.
%
%Auch der soziale Treffpunkt ist ins Internet umgezogen. Statt auf dem Marktplatz trifft man sich heute in Feeds, Chats und Foren.
%Wer wem begegnet und was man zu sehen bekommt, entscheiden dort aber Algorithmen von Konzernen.
%
%Die Agora gibt es also noch immer. Geändert haben sich die Regeln, nach denen man darin erfolgreich ist, und vor allem, wer sie
%schreibt.
%}
%
%
%\folie{Gesellschaftsordnung im Wandel}{%
%Heute bewegt sich unser Markt in einem spürbaren Spannungsfeld. Einerseits verlagert sich der Alltag immer weiter ins Digitale: Laut Bundesbank wurden erstmals 55\,\% aller Einkäufe bargeldlos bezahlt, während allein Plattformen wie Temu fast 21 Millionen Menschen pro Monat in Deutschland erreichen. Andererseits herrscht im Land eine deutliche Zurückhaltung – das Konsumklima liegt Ende 2026 tief im Minus, und die Sparquote ist mit rund 19\,\% so hoch wie kaum sonst in Europa.
%
%Hinter diesen Zahlen steckt ein tiefes Gefühl der Verunsicherung. Viele Menschen blicken skeptisch in den eigenen Geldbeutel, halten ihr Geld aus Vorsicht zusammen und haben das Vertrauen in politische Institutionen und soziale Medien weitgehend verloren. Man spürt, dass sich die Spielregeln verändern, ohne dass klar ist, wer sie eigentlich noch bestimmt.
%
%Dabei zeigt ein Blick zurück, wofür ein Markt im Kern steht: Er ist ein Spiegelbild der Gesellschaft. Die antike \textbf{Agora} war nicht nur ein Handelsplatz, sondern auch politischer Treffpunkt und sozialer Mittelpunkt. Hier kamen alle zusammen – Kaufleute, Gaukler und Richter. Der Markt versorgte die Menschen mit Lebensmitteln, und gleichzeitig nahmen Politik und Institutionen über ihn Einfluss auf das Zusammenleben.
%
%Damit diese Ordnung funktioniert, braucht es ein System, das auf Geben und Nehmen beruht. Wer etwas bekommen will, muss etwas einbringen, sei es Ware, Arbeit, Geld oder Vertrauen. Wer in diesem System zurechtkommen will, muss seine Regeln kennen – und darauf vertrauen können, dass sie für alle gelten.
%}
%

\folie{Zeitgeist}{%

Markt und Handel sind weit mehr als der bloße Austausch von Waren. Im Kern bilden sie ein gesellschaftliches Ökosystem, in dem Güter, Arbeit und Informationen zirkulieren. Sichtbar wird dieses Prinzip bereits in der antiken \textbf{Agora}, die Handelsplatz, politischer Treffpunkt und sozialer Mittelpunkt zugleich war. Hier kamen alle zusammen, Kaufleute, Gaukler und Richter. Während der Markt die Menschen mit Lebensmitteln versorgte, verhandelten Politik und Institutionen dort die Regeln des Zusammenlebens. Getragen wird ein solches System von einem Gleichgewicht aus Geben und Nehmen. Wer etwas erhalten will, muss etwas einbringen und die Spielregeln kennen. Handel ist damit vor allem ein verbindliches Regelwerk für alle Beteiligten.

Heute hat sich dieses Prinzip tiefgreifend gewandelt und ins Digitale verlagert. Informationen fließen über Plattformen wie Instagram, Facebook oder X, und auf Amazon, Temu oder Shein ist der Marktplatz rund um die Uhr geöffnet. Allein in Deutschland wurden 2024 täglich rund 400.000 Pakete von Temu und Shein zugestellt. Gingen die Rahmenbedingungen früher von sichtbaren Gemeinschaften und Institutionen aus, bestimmen heute unsichtbare Algorithmen und Konzerne, was überhaupt wahrgenommen wird und wer zu welchen Bedingungen teilnehmen darf. Bezahlt wird dabei längst nicht mehr nur mit Geld, sondern ebenso mit Daten und Aufmerksamkeit.

}

\folie{Zeitgeist -- Gefühle}{%

Gegenüber den heutigen Markt- und Handelsstrukturen wachsen bei vielen Menschen Wut und Frustration. Immer mehr gehen auf die Straße, weil in dieser Ordnung die Schere zwischen Arm und Reich immer größer wird. So besitzen die zehn reichsten Familien viermal so viel wie die untere Hälfte der Bevölkerung. Ein großer Teil dieses Vermögens steckt in Konzernen, denen auch die Plattformen gehören, über die wir uns informieren und austauschen. Damit verändert diese Entwicklung auch unser soziales Bewusstsein und die Art, wie wir miteinander kommunizieren. Wer heute auf dem Laufenden bleiben will, kommt an sozialen Netzwerken kaum noch vorbei. Für Jugendliche sind sie längst zur Grundvoraussetzung gesellschaftlicher Teilhabe geworden.

Gleichzeitig, Wissen ist frei zugänglich, Produkte sind nahezu grenzenlos verfügbar, und fast alles lässt sich rund um die Uhr mit einem Klick erreichen. Noch nie war es so einfach, sich eigenständig weiterzubilden, Preise und Qualität weltweit zu vergleichen oder mit geringem Startkapital selbst als Anbieter aufzutreten. In der Menschheitsgeschichte ist das einzigartig, wird im Alltag aber oft kaum noch bewusst wahrgenommen. Wer die dahinterliegenden Mechanismen und Algorithmen jedoch durchschaut, kann sie gezielt für sich nutzen und sich souverän darin behaupten.

Unsere heutige Wirtschafts- und Informationsordnung ist komplex, vielschichtig und greift tief in den Alltag ein. Da sie sowohl große Chancen als auch erhebliche Risiken birgt, kommt es darauf an, die Spielregeln des Handels zu verstehen, um die Vorteile auf persönlicher wie gesellschaftlicher Ebene zu nutzen und mit den Schattenseiten bewusst umzugehen.

}

% ================= Systemkarte: Kernelemente & Schalenmodell =================
\kartenfolie{Systemkarte}{%
  \begin{scope}[shift={(17.6,-12.0)}]
    \skKernfalse
    % ============================================================================
% Agora-Systemkarte (Kern-Trias + Schalenelemente), farbig
% Einbinden INNERHALB einer tikzpicture (Koordinaten um (0,0) zentriert).
% Benötigt: arrows.meta, calc. Farbe: skC (= hkared, einheitlich rot);
%           \newif\ifskKern im Preamble.
% Vor dem \input \skKerntrue setzen, um nur die Kernelemente zu zeichnen.
% ============================================================================
\providecolor{hkared}{RGB}{198, 42, 18}
\colorlet{skC}{hkared}  % einheitliche Farbe der Systemkarte
\providecommand{\pTag}{\textbf{[+]}}
\providecommand{\mTag}{\textbf{[--]}}
\providecommand{\pmTag}{\textbf{[$\pm$]}}

\tikzset{
  skCore/.style = {
    rounded corners=5pt, draw=#1!85!black, line width=1.4pt, fill=#1!10, text=black!85,
    minimum width=4.15cm, align=center, inner xsep=9pt, inner ysep=6pt
  },
  skShell/.style = {
    rounded corners=4pt, draw=#1!85!black, line width=1.15pt, fill=#1!5, text=black!85,
    minimum width=3.85cm, align=center, inner xsep=8pt, inner ysep=5.5pt
  },
  skAgora/.style = {
    ellipse, draw=skC!85!black, line width=2.0pt, fill=skC!12, text=black!85,
    minimum width=3.20cm, minimum height=1.35cm, align=center, font=\Large\bfseries
  },
  % Pfeilfarbe = Element, von dem die Wirkung ausgeht
  skArrow/.style = {
    line width=2.2pt, draw=#1!80!black, rounded corners=9pt, line cap=round, line join=round,
    -{Latex[round, length=3.4mm, width=3.2mm]},
    preaction={draw=white, line width=5.5pt, arrows=-}
  },
  skLbl/.style = {
    font=\fontsize{6.4}{7.3}\selectfont, text=#1!70!black, fill=white,
    inner xsep=2.4pt, inner ysep=1.2pt, align=center
  },
  skLblB/.style = {
    font=\fontsize{6.6}{7.5}\selectfont\bfseries, text=#1!70!black, fill=white,
    inner xsep=2.6pt, inner ysep=1.3pt, align=center
  },
}

% ---------------------------------------------------------------- Knoten
\node[skAgora] (AG) at (0.0, -0.25) {Agora};
\draw[skC!60, line width=0.65pt] (AG) ellipse (1.48cm and 0.60cm);

\node[skCore=skC] (ST) at (0.0, 3.50) {
  {\normalsize\bfseries Sozialer Treffpunkt}\\[2.5pt]
  \begin{minipage}{3.2cm}\scriptsize
    $\bullet$ Menschen\\ $\bullet$ Kultur\\ $\bullet$ Freizeit
  \end{minipage}};
\node[skCore=skC] (H) at (-5.6, -2.40) {
  {\normalsize\bfseries Handelsplatz}\\[2.5pt]
  \begin{minipage}{3.2cm}\scriptsize
    $\bullet$ Finanzsystem\\ $\bullet$ Digitalität\\ $\bullet$ Ressourcen
  \end{minipage}};
\node[skCore=skC] (P) at (5.6, -2.40) {
  {\normalsize\bfseries Politik}\\[2.5pt]
  \begin{minipage}{3.2cm}\scriptsize
    $\bullet$ Exekutive\\ $\bullet$ Legislative\\ $\bullet$ Judikative
  \end{minipage}};

\ifskKern\else
\node[skShell=skC] (I) at (-11.3, 3.15) {
  {\normalsize\bfseries Infrastruktur}\\[2.5pt]
  \scriptsize $\bullet$ Verkehr \enspace $\bullet$ Energie \enspace $\bullet$ Netze};
\node[skShell=skC] (N) at (11.3, 3.15) {
  {\normalsize\bfseries Normen \& Werte}\\[2.5pt]
  \scriptsize $\bullet$ Vertrauen \enspace $\bullet$ Solidarität \enspace $\bullet$ Fairness};
\node[skShell=skC] (U) at (0.0, -7.35) {
  {\normalsize\bfseries Unternehmen}\\[2.5pt]
  \scriptsize $\bullet$ Konzerne \enspace $\bullet$ Mittelstand \enspace $\bullet$ Start-ups};
\fi

% ---------------------------------------------------------------- Pfeile
\begin{scope}
  % --- Kern-Trias ---
  \draw[skArrow=skC]    ([xshift=-10mm]H.north) |- ([yshift=4mm]ST.west);
  \draw[skArrow=skC]   ([yshift=-6mm]ST.west)  -| ([xshift=10mm]H.north);
  \draw[skArrow=skC] ([xshift=10mm]P.north)  |- ([yshift=4mm]ST.east);
  \draw[skArrow=skC]   ([yshift=-6mm]ST.east)  -| ([xshift=-10mm]P.north);
  \draw[skArrow=skC]    ([yshift=5mm]H.east)    -- ([yshift=5mm]P.west);
  \draw[skArrow=skC] ([yshift=-5mm]P.west)   -- ([yshift=-5mm]H.east);

  \ifskKern\else
  % --- Kern <-> Schale ---
  \draw[skArrow=skC]   ([xshift=4mm]I.north)   |- (-6.0, 5.30) -| ([xshift=-10mm]ST.north);
  \draw[skArrow=skC]   ([xshift=10mm]ST.north) |- (6.0, 5.30)  -| ([xshift=-4mm]N.north);
  \draw[skArrow=skC]  ([xshift=4mm]N.north)   |- (6.0, 6.00)  -| ([xshift=4mm]ST.north);
  \draw[skArrow=skC]   ([xshift=6mm]I.south)   |- ([yshift=5mm]H.west);
  \draw[skArrow=skC]    ([yshift=-5mm]H.west)   -| ([xshift=-4mm]I.south);
  \draw[skArrow=skC]  ([xshift=-6mm]N.south)  |- ([yshift=5mm]P.east);
  \draw[skArrow=skC] ([yshift=-5mm]P.east)   -| ([xshift=4mm]N.south);
  \draw[skArrow=skC]    ([xshift=-8mm]H.south)  |- ([yshift=3mm]U.west);
  \draw[skArrow=skC]    ([xshift=-12mm]U.north) |- (-3.0, -5.00) -| ([xshift=10mm]H.south);
  \draw[skArrow=skC] ([xshift=8mm]P.south)   |- ([yshift=3mm]U.east);
  \draw[skArrow=skC]    ([xshift=12mm]U.north)  |- (3.0, -5.00)  -| ([xshift=-10mm]P.south);
  % --- Schale <-> Schale ---
  \draw[skArrow=skC]   ([xshift=-8mm]I.north)  |- (0.0, 6.80)   -| ([xshift=8mm]N.north);
  \draw[skArrow=skC]  ([xshift=12mm]N.south)  |- ([yshift=-3mm]U.east);
  \draw[skArrow=skC]    ([yshift=-3mm]U.west)   -| ([xshift=-12mm]I.south);
  \fi
\end{scope}

% ---------------------------------------------------------------- Beschriftungen
\begin{scope}
  \node[skLblB=skC]    at ($([xshift=-10mm]H.north)+(0,3.0)$) {bietet Waren, Jobs\\[-1pt]\& Güter \pTag};
  \node[skLblB=skC]   at ($([xshift=10mm]H.north)+(0,1.7)$)  {schafft Konsum\\[-1pt]\& Nachfrage \pTag};
  \node[skLblB=skC] at ($([xshift=10mm]P.north)+(0,3.0)$)  {schützt Rechte\\[-1pt]\& setzt Rahmen \pmTag};
  \node[skLblB=skC]   at ($([xshift=-10mm]P.north)+(0,1.7)$) {wählt \& artikuliert\\[-1pt]Anliegen \pTag};
  \node[skLblB=skC]    at ($(H.east)!0.5!(P.west)+(0,0.5)$)   {finanziert über Steuern \& Abgaben \pTag};
  \node[skLblB=skC] at ($(H.east)!0.5!(P.west)+(0,-0.5)$)  {setzt Marktregeln \& Gesetze \pmTag};

  \ifskKern\else
  \node[skLbl=skC]   at (-6.0, 5.30) {schafft Zugang, Mobilität \& Raum \pTag};
  \node[skLbl=skC]   at ($([xshift=6mm]I.south |- H.west)!0.5!(H.west)+(0,0.5)$)  {ermöglicht Logistik\\[-1pt]\& Lieferketten \pTag};
  \node[skLbl=skC]    at ($([xshift=-4mm]I.south |- H.west)!0.5!(H.west)+(0,-0.5)$) {treibt Ausbau \&\\[-1pt]Digitalisierung \pTag};
  \node[skLbl=skC]   at (4.6, 5.30)  {verhandelt \& prägt Leitbilder \pTag};
  \node[skLbl=skC]  at (7.6, 6.00)  {stiften Vertrauen \& Zusammenhalt \pTag};
  \node[skLbl=skC]  at ($([xshift=-6mm]N.south |- P.east)!0.5!(P.east)+(0,0.5)$) {setzen Maßstäbe\\[-1pt]für Gesetze \pTag};
  \node[skLbl=skC] at ($([xshift=4mm]N.south |- P.east)!0.5!(P.east)+(0,-0.5)$) {verankert Werte\\[-1pt]im Recht \pTag};
  \node[skLbl=skC]    at ($([xshift=-8mm]H.south)+(0,-1.9)$) {liefert Kapital,\\[-1pt]Arbeit \& Rohstoffe \pTag};
  \node[skLbl=skC]    at (-3.0, -5.00) {schaffen Güter,\\[-1pt]Dienste \& Wert \pTag};
  \node[skLbl=skC] at ($([xshift=8mm]P.south)+(0,-1.9)$)  {reguliert, besteuert\\[-1pt]\& fördert \pmTag};
  \node[skLbl=skC]    at (3.0, -5.00)  {betreiben Lobbyismus\\[-1pt]\& schaffen Jobs \pmTag};
  \node[skLblB=skC] at ($([xshift=12mm]N.south)+(0,-4.2)$)  {fordern Ethik, Fairness\\[-1pt]\& Nachhaltigkeit \mTag};
  \node[skLblB=skC]   at ($([xshift=-12mm]I.south)+(0,-4.2)$) {brauchen \& verschleißen\\[-1pt]Netze/Energie \mTag};
  \node[skLblB=skC]  at (0.0, 6.80) {vernetzt Gesellschaft \& beschleunigt Diskurs \pTag};
  \fi
\end{scope}

  \end{scope}
}


\folie{Lern- \& Kommunikationsziele}{%
  \begin{enumerate}
    \item \textbf{Reziprozität und Wechselwirkungen}\\
          Jede Handlung im Markt löst eine Gegenreaktion aus.
    \item \textbf{Reale Probleme des heutigen Handels}\\
          Ungleichheit, Ausbeutung und Plattformmacht werden spürbar.
    \item \textbf{Abhängigkeiten}\\
          Niemand handelt allein, alle sind auf Infrastruktur, Partner und Regeln angewiesen.
    \item \textbf{Vertrauen als Währung}\\
          Ein guter Ruf öffnet Türen, gebrochenes Vertrauen ist schwer zurückzugewinnen.
    \item \textbf{Wer die Regeln versteht, gewinnt}\\
          Wer Mechanismen und Machtverhältnisse durchschaut, kann sie gezielt nutzen.
    \item \textbf{Eigener Vorteil oder Gemeinwohl}\\
          Kurzfristiger Gewinn kann dem gesamten System schaden.
  \end{enumerate}
}

% ================= Quellen =================
\folie{Quellen}{%
\fontsize{10}{12.5}\selectfont\setlength{\parskip}{0.25em}\raggedright
\textbf{Grundlagen Handel, Geld \& Vertrauen}\\
Jevons, W. S. (1875). \emph{Money and the Mechanism of Exchange}. London, H. S. King.\\
Simmel, G. (1900). \emph{Philosophie des Geldes}. Leipzig, Duncker \& Humblot.\\
Luhmann, N. (1968). \emph{Vertrauen. Ein Mechanismus der Reduktion sozialer Komplexität}. Stuttgart, Enke.\\
Graeber, D. (2011). \emph{Debt: The First 5,000 Years}. New York, Melville House.

\textbf{Plattformen, Konsum \& Gamification}\\
Zuboff, S. (2019). \emph{The Age of Surveillance Capitalism}. New York, PublicAffairs.\\
Carrington, M. J.; Neville, B. A.; Whitwell, G. J. (2010). Why Ethical Consumers Don't Walk Their Talk. \emph{Journal of Business Ethics} 97, S.~139 bis 158.\\
BEUC (2024). \emph{Taming Temu}. Beschwerde europäischer Verbraucherverbände zu manipulativen Designmustern. \href{https://www.beuc.eu}{www.beuc.eu}\\
Handelsverband Deutschland (HDE) (2024). Pressemitteilungen zum Paketvolumen von Temu und Shein. \href{https://einzelhandel.de}{einzelhandel.de}

\textbf{Zahlungsverhalten}\\
Deutsche Bundesbank (2026). \emph{Zahlungsverhalten in Deutschland 2025}. \href{https://www.bundesbank.de}{www.bundesbank.de}

\textbf{Lieferketten \& Rotes Meer}\\
UNCTAD (2024). \emph{Review of Maritime Transport 2024}. Genf, United Nations. \href{https://unctad.org/rmt2024}{unctad.org/rmt2024}\\
Drewry (2024 bis 2026). \emph{World Container Index}. \href{https://www.drewry.co.uk}{www.drewry.co.uk}\\
Shanghai Shipping Exchange. \emph{Shanghai Containerized Freight Index (SCFI)}. \href{https://en.sse.net.cn}{en.sse.net.cn}

\textbf{Regulierung}\\
Lieferkettensorgfaltspflichtengesetz (LkSG) vom 16.07.2021, BGBl. I S.~2959.\\
Richtlinie (EU) 2024/1760 über die Sorgfaltspflichten von Unternehmen im Hinblick auf Nachhaltigkeit (CSDDD), ABl. L, 05.07.2024.
}

\end{document}
